\documentclass[grey]{hipstercv}
% available options are: darkhipster, lighthipster, pastel, allblack, grey, verylight
\usepackage[utf8]{inputenc}
\usepackage[default]{raleway}
\usepackage[margin=1cm, a4paper]{geometry}


%------------------------------------------------------------------ Variablen

\newlength{\rightcolwidth}
\newlength{\leftcolwidth}
\setlength{\leftcolwidth}{0.3\textwidth}
\setlength{\rightcolwidth}{0.65\textwidth}

%------------------------------------------------------------------
\title{CV -- Arthur}
\author{Séptima Ola}
\date{Agosto 2026}

\pagestyle{empty}
\begin{document}


\thispagestyle{empty}
%-------------------------------------------------------------

\section*{Start}



\header{\bg{headerfontbox}{headerfontboxfont}{Bajista} }{\bgupper{headerfontbox}{headerfontboxfont}{\bfseries\Huge Arthur}}{\bg{headerfontbox}{headerfontboxfont}{\large Séptima Ola}}{arthur.png}{headerblue}{3.5cm}{2cm}



%------------------------------------------------

% hier muss die "unsichtbare" Überschrift rein, weil er sonst nicht die Paracols startet... komisch...
\subsection*{}
\vspace{4em}

\setlength{\columnsep}{1.5cm}
\columnratio{0.3}[0.65]
\begin{paracol}{2}
\hbadness5000
%\backgroundcolor{c[1]}[rgb]{1,1,0.8} % cream yellow for column-1 %\backgroundcolor{g}[rgb]{0.8,1,1} % \backgroundcolor{l}[rgb]{0,0,0.7} % dark blue for left margin

\paracolbackgroundoptions

% 0.9,0.9,0.9 -- 0.8,0.8,0.8


\footnotesize
{\setasidefontcolour
\bgupper{cvgreen}{white}{Perfil} \\
\bg{cvgreen}{white}{Datos personales} \\

\begin{tabular}{ll}
\faMale&Arthur \\
\faGlobe& Edo. de México, MX \\
\faBirthdayCake&Septiembre 1992 \\
\faMapMarker&Estado de México \\
\end{tabular}

\bigskip

\bg{cvgreen}{white}{Especialidad} \\

Bajo eléctrico ~•~ Salsa ~•~ Reggae ~•~ Ska

\bigskip

\bgupper{cvgreen}{white}{Habilidades} \\
\bg{cvgreen}{white}{Idiomas}
\bigskip


\begin{minipage}[t]{\leftcolwidth}
\begin{tabular}{l | ll}
\textbf{Español} & C2 & {\phantom{x}\footnotesize lengua materna} \\
\textbf{Inglés} & A2 & \pictofraction{\faCircle}{cvpurple}{1}{black!30}{3}{\tiny}
\end{tabular}
\end{minipage}

\bigskip

\bg{cvgreen}{white}{Intereses}\\

\bubblediagram{el bajo, \textbf{el groove}, salsa, reggae, ska, rocksteady, \textbf{coleccionar}\\ \textbf{bajos}, probar cuerdas, ensayar, \textbf{música}\\ \textbf{latina}}

\bigskip

\bg{cvgreen}{white}{Estilos y géneros} \\

\begin{minipage}[t]{0.3\textwidth}
\begin{tabular}{r @{\hspace{0.5em}}l}
     \bg{skilllabelcolour}{iconcolour}{Salsa} &  \barrule{0.9}{0.5em}{cvpurple}\\
     \bg{skilllabelcolour}{iconcolour}{Reggae} & \barrule{0.9}{0.5em}{cvgreen} \\
     \bg{skilllabelcolour}{iconcolour}{Ska} & \barrule{0.8}{0.5em}{cvpurple} \\
     \bg{skilllabelcolour}{iconcolour}{Rocksteady} & \barrule{0.7}{0.5em}{cvpurple} \\
     \bg{skilllabelcolour}{iconcolour}{Versátil} & \barrule{0.75}{0.5em}{cvpurple} \\
\end{tabular}
\end{minipage}


\bigskip

\scalebox{0.8}{
\iconcross{\Huge}{white}{cvred}{\color{black!30}\faBook}{\href{mailto:septimaola@gmail.com}{\faEnvelopeO}}{\faPhone}{\faMusic}
}

\phantom{turn the page}

\phantom{turn the page}
}
%-----------------------------------------------------------
\switchcolumn

\small
\section*{Trayectoria}

\begin{tabular}{r| p{0.4\textwidth} c}
    \cvevent{2024--presente}{Séptima Ola}{Bajo eléctrico}{Edo. de México \color{cvred}}{Bajo que camina entre salsa y reggae, dando identidad rítmica al sonido de la banda.}{music_logo.png} \\
    \cvevent{2018--2023}{El Croar Norteño}{Bajo eléctrico}{Edo. de México \color{cvred}}{Proyecto de música versátil donde afianzó técnica y criterio en el bajo.}{music_logo.png} \\
    \cvevent{2015--2018}{Arthur en sus días}{Bajo eléctrico}{Edo. de México \color{cvred}}{Proyecto versátil junto con Alfred. Base del sonido que hoy define su estilo.}{music_logo.png}
\end{tabular}

\vspace{4em}

\begin{minipage}[t]{0.4\textwidth}
\section*{Formación}
\begin{tabular}{r p{0.6\textwidth} c}
    \cvdegree{2010--presente}{Autodidacta}{Bajo eléctrico}{Trayectoria empírica \color{headerblue}}{Horas con el bajo en mano, copiando líneas de salsa y reggae hasta hacerlas propias.}{music_logo.png} \\
    \cvdegree{2015}{Taller de bajo}{Curso}{Escena local \color{headerblue}}{Complemento ocasional; la base siempre fue el aprendizaje autodidacta.}{music_logo.png}
\end{tabular}
\end{minipage}\hfill
\begin{minipage}[t]{0.16\textwidth}
\section*{Pasatiempos}
\hobbyicon{\color{iconcolour}\faMusic}{Bajo}{cvgreen}{\iconsize}{2em} \hfill
\hobbyicon{\color{iconcolour}\faHeadphones}{Discos}{cvorange}{\iconsize}{2em}

\hobbyicon{\color{iconcolour}\faUsers}{Ensayar}{cvpurple}{\iconsize}{2em} \hspace{1em}
\hobbyicon{\color{iconcolour}\faWrench}{Cuerdas}{headerblue}{\iconsize}{2em}
\end{minipage}

\vspace{4em}

\begin{minipage}[t]{0.35\textwidth}
\section*{Influencias}
\begin{tabular}{>{\footnotesize\bfseries}r >{\footnotesize}p{0.55\textwidth}}
    & Salsa clásica (bajistas) \\
    & Bob Marley \& The Wailers \\
    & Reggae roots jamaiquino \\
    & Willie Colón \\
    & Rocksteady jamaiquino
\end{tabular}

\section*{Fortalezas}
\cvtag{autodidacta}
\cvtag{groove}
\cvtag{disciplina}
\cvtag{versatilidad}

\section*{Referencias}
\cvkeyword{Alfred (Séptima Ola)}{cvgreen}{iconcolour}
\cvkeyword{El Croar Norteño}{cvgreen}{iconcolour}
\end{minipage}\hfill
\begin{minipage}[t]{0.3\textwidth}
\section*{Equipo}
\begin{tabular}{>{\footnotesize\bfseries}r >{\footnotesize}p{0.65\textwidth}}
    & Colección de bajos eléctricos \\
    & Variedad de cuerdas (cada instrumento inspira un color distinto)
\end{tabular}

\section*{Frase}
\begin{quote}\footnotesize
``Creo en la música como un puente entre almas; el bajo es el cable que une el ritmo con el corazón.''
\end{quote}
\end{minipage}









\vfill{} % Whitespace before final footer

%----------------------------------------------------------------------------------------
%	FINAL FOOTER
%----------------------------------------------------------------------------------------
\setlength{\parindent}{0pt}
\begin{minipage}[t]{\rightcolwidth}
\begin{center}\fontfamily{\sfdefault}\selectfont \color{black!70}
{\small Arthur \icon{\faEnvelopeO}{cvpurple}{} Séptima Ola \icon{\faMapMarker}{cvpurple}{} Edo. de México \icon{\faPhone}{cvpurple}{} TBD \newline\icon{\faAt}{cvpurple}{} \protect\url{septimaola@gmail.com}
}
\end{center}
\end{minipage}


\end{paracol}

\end{document}